\documentclass[10pt,a4paper]{article}
\usepackage{fontspec,amsmath,amssymb,graphicx}
\usepackage[margin=18mm]{geometry}
\setmainfont{Times New Roman}
\setlength{\parindent}{0pt}
\setlength{\parskip}{5pt}
\allowdisplaybreaks
\sloppy
\emergencystretch=2em
\begin{document}
\section*{BT\_3\_1\_3\_08 — Lời giải bổ sung}
\subsection*{Kiến thức và trình tự giải}
\textbf{Bước 1 — Đọc đúng điều kiện biên.} Biên $u(0,t)=0$ dùng mở rộng lẻ; biên $u_x(0,t)=0$ dùng mở rộng chẵn. Với dữ liệu trên $x\ge0$, đặt $\widetilde f(-s)=-f(s)$ hoặc $f(s)$ tương ứng, làm tương tự cho $g$. Mở rộng lẻ là $C^2$ nếu $f(0)=f''(0)=0$, và mở rộng lẻ của $g$ là $C^1$ nếu $g(0)=0$. Mở rộng chẵn đòi $f'(0)=g'(0)=0$ cho mức trơn ấy.

\textbf{Bước 2 — Lập nguyên hàm sau khi phản xạ.} Đặt $H(s)=\int_0^s\widetilde g(r)dr$. Nếu $\widetilde g$ lẻ thì $H$ chẵn; nếu $\widetilde g$ chẵn thì $H$ lẻ. Chính nguyên hàm, chứ không phải vận tốc, xuất hiện trong


\[
u=\frac{\widetilde f(x-at)+\widetilde f(x+at)}2+\frac{H(x+at)-H(x-at)}{2a}.
\]


\textbf{Bước 3 — Tách miền trước và sau khi chạm biên.} Với $u(0,t)=h(t)$, ở $x\ge at$ dùng d'Alembert với dữ liệu gốc. Ở $0<x<at$:


\[
u=h(t-x/a)+\frac{f(x+at)-f(at-x)}2+\frac1{2a}\int_{at-x}^{at+x}g(s)ds.
\]


Với $u_x(0,t)=h(t)$, dùng nghiệm của các mở rộng chẵn rồi trừ $a\int_0^{t-x/a}h(s)ds$ trên miền $x<at$. Tính riêng đạo hàm theo $x$ của hạng này để kiểm tra biên.

\textbf{Bước 4 — Trên đoạn hữu hạn, phản xạ ở cả hai đầu.} Dirichlet ở hai đầu cho mở rộng lẻ tại $0,L$, chu kỳ không gian $2L$. Neumann ở hai đầu cho mở rộng chẵn tại $0,L$, cũng chu kỳ $2L$. Một đầu lẻ, một đầu chẵn cho $\widetilde f(s+2L)=-\widetilde f(s)$ và chu kỳ $4L$, tương tự $g$. Vì thế chu kỳ thời gian thường là $2L/a$ hoặc $4L/a$. Với Neumann hai đầu còn phải xét trung bình $g$: trung bình khác không sinh hạng tuyến tính theo $t$, nên nghiệm không tuần hoàn.

\textbf{Bước 5 — Kiểm tra góc và đường ghép.} Biên Dirichlet không thuần nhất cần $h(0)=f(0)$, $h'(0)=g(0)$, $h''(0)=a^2f''(0)$ để có nghiệm $C^2$ cho phương trình thuần nhất. So sánh hai vết trên $x=at$, gồm $u,u_x,u_t$ rồi đạo hàm cấp hai. Nếu không tương thích, nêu rõ mức trơn còn lại thay vì khẳng định nghiệm cổ điển trên toàn miền.
\subsection*{Lời giải áp dụng cho bài}
\textbf{Bước 1. Dùng công thức Dirichlet với thời điểm trễ.} Ở $x\ge at$ sóng chưa chạm biên; ở $x<at$ dùng $h(t-x/a)$ và phản xạ lẻ. Dưới đây giữ hai miền khi biên làm thay đổi vận tốc.

(i) $a=2$, $f=\sin x$, $g=\cos2x$, $h=0$. Mở rộng $g$ lẻ, $H(s)=\sin(2|s|)/2$. Do đó
\[u=\sin x\cos2t+\begin{cases}\cos2x\sin4t/4&x\ge2t,\\\sin2x\cos4t/4&0<x<2t.\end{cases}\]
$g(0)=1\ne h'(0)$ nên không $C^1$ trên mặt sóng. Sau khi mặt sóng đi qua mỗi $x$, công thức tuần hoàn chu kỳ $\pi$; toàn nghiệm không có chu kỳ chung vì hai nhánh khác nhau.

(ii) $a=3$, $f=g=0$, $h=t/(1+t)$. Đặt $\tau=t-x/3$: $u=0$ khi $\tau\le0$, $u=\tau/(1+\tau)$ khi $\tau>0$. Trên tia $x=ct$, giới hạn là $1$ nếu $0\le c<3$, và $0$ nếu $c\ge3$. Nghiệm không tuần hoàn: tại mỗi $x$ cố định, sau khi mặt sóng đến, $u$ tăng nghiêm ngặt về $1$.

(iii) $a=1$, mở rộng lẻ $f$ là $s|s|$, còn $g=\sin s$ đã lẻ:
\[u=\tfrac12[(x+t)|x+t|+(x-t)|x-t|]+\sin x\sin t.\]
Nghiệm $C^1$, không $C^2$ trên $x=t$ vì $f''(0)=2\ne0$, và không tuần hoàn.

(iv) $a=5$, $f=0$, $g=\sin x-\sin3x/3$ có $H=-\cos s+\cos3s/9$. Suy $u=\sin x\sin5t/5-\sin3x\sin15t/45$, trơn, chu kỳ $2\pi/5$.

(v) $a=2$, $f=\cos x-\sin x$, $g=0$, $h=\cos2t-\sin2t$. Ngoài mặt sóng, half-sum của $f$ là $(\cos x-\sin x)\cos2t$. Trong mặt sóng:
\[h(t-x/2)+\tfrac12[f(x+2t)-f(2t-x)]=\cos x(\cos2t-\sin2t).\]
Vậy
\[u=\begin{cases}(\cos x-\sin x)\cos2t&x\ge2t,\\\cos x(\cos2t-\sin2t)&0<x<2t.\end{cases}\]
Công thức cũ $\cos(x-2t)-\sin(x+2t)$ không thỏa $u_t(x,0)=0$. Nghiệm đúng liên tục, không $C^1$ vì $h'(0)=-2\ne0$. Chỉ nhánh sau mặt sóng có chu kỳ $\pi$ ở từng điểm.

(vi) $a=1$, $g=\chi_{[1,2]}-\chi_{[3,4]}$. Mở rộng lẻ có nguyên hàm chẵn $H(s)=[|s|-1]_0^1-[|s|-3]_0^1$, nên $u=(H(x+t)-H(x-t))/2$. Tại $x$ cố định, hai giá trị $H$ đều bằng không sau $t>x+4$, nên $u=0$; đây là xung quá độ, không tuần hoàn. Nghiệm liên tục nhưng có gãy cấp một khi các cận gặp $\pm1,\pm2,\pm3,\pm4$.

\textbf{Bước 2. Kiểm tra chung.} Thế $x=0$ vào nhánh trong cho đúng $h$; thế $t=0$ vào nhánh ngoài cho $f,g$. So hai vết trên $x=at$ theo ba điều kiện tương thích, không dùng tính tuần hoàn của nhánh trong để kết luận tuần hoàn toàn miền.

\textbf{Chi tiết hai nhánh vận tốc ở (i).} Ngoài mặt sóng,
\[\frac14\int_{x-2t}^{x+2t}\cos2s\,ds
=\frac{\sin(2x+4t)-\sin(2x-4t)}8=\frac14\cos2x\sin4t.\]
Trong mặt sóng, cận dưới là $2t-x$, nên hiệu sine bằng $\sin(4t+2x)-\sin(4t-2x)=2\cos4t\sin2x$; chia $8$ cho nhánh đã ghi.

\textbf{Chi tiết rút gọn (v).} Trong miền phản xạ, biên là $\cos(2t-x)-\sin(2t-x)$. Half-difference của $f$ bằng
\[-\sin2t\sin x-\cos2t\sin x.\]
Khai triển biên thành $\cos2t\cos x+\sin2t\sin x-\sin2t\cos x+\cos2t\sin x$ rồi cộng, hai hạng chứa $\sin x$ triệt tiêu, còn $\cos x(\cos2t-\sin2t)$.
\begin{center}\includegraphics[width=.97\linewidth,height=11cm,keepaspectratio]{Do_thi/BT_3_1_3_08.png}\end{center}
\textit{Đồ thị minh họa; với nghiệm số, xem phương pháp và tham số trong mã Python kèm theo.}
\begin{center}\includegraphics[width=.97\linewidth,height=11cm,keepaspectratio]{Do_thi/BT_3_1_3_08_characteristics.png}\end{center}\end{document}